The assessment followed a structured Vulnerability Assessment and Penetration
Testing (VAPT) methodology. The methodology was designed to move from
initial planning and reconnaissance to vulnerability identification,
validation, documentation, and remediation.

\begin{figure}[!htb]
\centering
\begin{tikzpicture}[
    node distance=0.25cm,
    every node/.style={
        font=\small\bfseries,
        align=center
    },
    phase/.style={
        rectangle,
        rounded corners,
        draw=primary,
        fill=lightblue,
        minimum width=0.82\textwidth,
        minimum height=0.65cm
    },
    arrow/.style={
        ->,
        thick,
        color=secondary
    }
]

\node[phase] (p1) {Phase 1\\Planning and Scoping};
\node[phase, below=of p1] (p2) {Phase 2\\Reconnaissance};
\node[phase, below=of p2] (p3) {Phase 3\\Vulnerability Assessment};
\node[phase, below=of p3] (p4) {Phase 4\\Manual Verification and Proof of Concept};
\node[phase, below=of p4] (p5) {Phase 5\\Post-Exploitation and Impact Analysis};
\node[phase, below=of p5] (p6) {Phase 6\\Reporting and Risk Assessment};
\node[phase, below=of p6] (p7) {Phase 7\\Remediation and Retesting};

\draw[arrow] (p1) -- (p2);
\draw[arrow] (p2) -- (p3);
\draw[arrow] (p3) -- (p4);
\draw[arrow] (p4) -- (p5);
\draw[arrow] (p5) -- (p6);
\draw[arrow] (p6) -- (p7);

\end{tikzpicture}
\caption{VAPT assessment methodology}
\end{figure}

Each phase shown above is documented as its own chapter later in this
report: Phase 1 in Chapter~\ref{chap:phase1}, Phase 2 in
Chapter~\ref{chap:phase2}, Phase 3 in Chapter~\ref{chap:phase3},
Phase 4 in Chapter~\ref{chap:phase4}, Phase 5 in
Chapter~\ref{chap:phase5}, Phase 6 in Chapter~\ref{chap:phase6}, and
Phase 7 in Chapter~\ref{chap:phase7}.